\section{Simulated Annealing Results}
\label{app:sa}

Starting from $T(7,1,9)$ ($|V|=78$), our simulated annealing
search discovered a smaller Leontovich graph with $|V|=76$ via
edge mutation and vertex contraction.  The result was verified
independently with exact integer arithmetic
(\texttt{scripts/verify\_76.py}).

The 76-vertex graph is a tree with 75 edges and degree sequence
\[
  [10, 10, 10, 10, 10, 9, 9, 7, \underbrace{2, \ldots, 2}_{7},
   \underbrace{1, \ldots, 1}_{61}].
\]
It retains the star-of-stars topology of $T(7,1,9)$ but with
asymmetric leaf counts at the bridge layer (two branches each lost one
leaf).

\begin{table}[H]
\centering
\caption{Active quotient spectra along the pruning sequence. Asymmetric pruning breaks identical-branch symmetry, expanding the quotient matrix and lifting new active eigenvalues $\lambda_2^*$ that govern the Leontovich crossover.}
\label{tab:sa_hits}
\begin{tabular}{@{}clccccc@{}}
  \toprule
  $|V|$ & Structure & Crossover $n$ & $d$ & $\lambda_2^*/\lambda_1$ & $\lambda_1$ & $\lambda_2^*$ \\
  \midrule
  78 & $T(7,1,[9^7])$     & 13  & 2 & 0.6877 & 3.3973 & 2.3364 \\
  77 & $T(7,1,[9^6,8^1])$ & 15  & 2 & 0.8911 & 3.3874 & 3.0187 \\
  76 & $T(7,1,[9^5,8^2])$ & 17  & 2 & 0.8997 & 3.3766 & 3.0380 \\
  \midrule
  75 & $T(7,1,[9^4,8^3])$ & --- & --- & 0.9089 & 3.3647 & 3.0581 \\
  \bottomrule
\end{tabular}
\end{table}

\noindent
The 76-vertex graph has an active secondary quotient eigenvalue of $\lambda_2^* \approx 3.0380$. Its first crossover
at $n = 17$ with $d = 2$ was verified by exact integer computation
of $\Hom(P_n, H) - \Hom(E_n^{(2)}, H) = 1{,}993{,}646 > 0$.
For $n \ge 17$, the sign of $\Delta_n$ strictly alternates between
even and odd $n$, indicating destructive interference from the
active secondary eigenvalue.

A 20-seed swarm of $5 \times 10^6$ SA steps each (totalling
$10^8$ evaluations) was run from $T(7,1,9)$ at temperature 3.0.
Every seed converged to a 76-vertex graph, with an intermediate
$|V|=77$ Leontovich graph consistently discovered along the path
(typically within the first 10 steps). No seed found a graph with
$|V| < 76$.  This strongly suggests that 76 is the global minimum
within the star-of-stars basin, and that the descent
$78 \to 77 \to 76$ represents the unique pruning sequence.

\begin{figure}[H]
\centering
\includegraphics[width=0.75\textwidth]{figures/leontovich_76.pdf}
  \caption{The 76-vertex Leontovich graph $T(7,1,[9^5,8^2])$. Red: root
  (degree~7). Blue: bridges (degree~2). Orange: hubs (degree~9--10).
  Green: leaves (degree~1). The asymmetric pruning of one leaf from
  each of two branches in $T(7,1,9)$ expands the active quotient
  spectrum, shifting the first crossover from $n=13$ to $n=17$.}
\label{fig:leo76}
\end{figure}

\subsection*{Symmetry breaking and the active quotient spectrum}

The pruning sequence admits a clean analytic
description that resolves the crossover behaviour. In the fully symmetric
graph $T(7,1,[9^7])$, the adjacency matrix possesses highly degenerate
eigenvalues at $\lambda = \sqrt{10}$ (arising from the isolated 9-leaf hubs).
However, the eigenvectors associated with this eigenspace are antisymmetric
across the identical branches. Consequently, they are perfectly orthogonal
to the all-ones vector $\mathbf{1}$ and the initial degree vector, meaning
their projection coefficients in the transfer matrix expansion
$\Hom(G_n, H) = \mathbf{1}^T A^{n-4} \mathbf{h}_0$ are identically zero.
These orthogonal eigenvectors contribute nothing to the homomorphism counts.

Instead, the crossover dynamics are governed exclusively by the \emph{active}
spectrum---the eigenvalues of the quotient matrix over the graph's
automorphism orbits. For the fully symmetric graph $T(7,1,[9^7])$, the
active secondary eigenvalue is $\lambda_2^* \approx 2.3364$.

When leaves are pruned asymmetrically, the perfect identical-branch symmetry
is broken. For $T(7,1,[9^k, 8^{7-k}])$, the orbits split and the quotient
matrix expands from $4 \times 4$ to $7 \times 7$. This symmetry breaking
activates a new secondary eigenvalue $\lambda_2^*$ that directly interferes
with the Perron mode.

The Leontovich crossover arises from destructive interference between the
dominant Perron mode $\lambda_1^n$ and the secondary active mode $(\lambda_2^*)^n$
in the difference $\Hom(P_n, H) - \Hom(E_n^{(d)}, H)$. As branches are pruned,
$\lambda_1$ decreases while the active secondary eigenvalue $\lambda_2^*$
increases (see Table~\ref{tab:sa_hits}), narrowing the spectral gap.
At $|V| = 76$, the interference from $\lambda_2^* \approx 3.0380$ is just
sufficient to drive the ratio $\Hom(E_n^{(2)}, H)/\Hom(P_n, H)$ below~1 at
$n = 17$. The local SMT experiment for $|V| = 75$ times out at its first
depth-2 threshold, $n=13$, so it does not establish minimality within this
pruning family.

This yields a clean dichotomy: \emph{within the star-of-stars basin,
asymmetric pruning breaks orbital symmetry and lifts new active eigenvalues
into the quotient spectrum.} The Leontovich property is lost when this
symmetry-breaking pushes the active secondary eigenvalue too close to the
spectral radius. The 76-vertex graph is the smallest witness found by the
search, not a certified minimum; Section~\ref{sec:local-pruning-audit}
describes the incomplete local audit.

\subsection*{Deeper symmetric trees}

An overflow-safe sweep (\texttt{\_\_int128} exact arithmetic with
\texttt{\_\_builtin\_mul\_overflow}) over 596,891 branching sequences
$T(d_1, \ldots, d_k)$ with branching depth $k \le 5$ (at most six
orbits) and $|V| \le 2{,}000$
found 7,016 Leontovich instances and confirms that $T(7,1,9)$ at
$|V| = 78$ is the vertex-count minimum within the swept symmetric
branching-sequence family (depth $k\le5$, $|V|\le2{,}000$). For each
threshold~$n$, the smallest 5-orbit tree is
consistently \emph{larger} than the corresponding 4-orbit tree:
for example, at $n = 13$, $T(1,8,1,12)$ ($|V| = 114$) versus
$T(7,1,9)$ ($|V| = 78$).  Deeper trees do extend the reachable
threshold range---5-orbit entries reach $n \ge 61$ where no 4-orbit
tree appears---but at vertex counts exceeding 500.

An earlier floating-point sweep produced false positives (e.g.,
$T(1,1,1,21)$ at $|V| = 25$ was incorrectly flagged as Leontovich).
Exact integer verification showed $\Hom(E_n, H) > \Hom(P_n, H)$ at
all~$n$ for these candidates.  The root cause was signed
\texttt{\_\_int128} overflow (undefined behaviour under \texttt{-O3}),
which silently corrupted the comparisons; the
\texttt{\_\_builtin\_mul\_overflow} intrinsic eliminates this class
of error.

\subsection*{The Diophantine threshold frontier}

The minimum vertex counts in Table~\ref{tab:landscape_full} are
explained algebraically and certified computationally in the tested range.

The audited frontier is restricted to the subfamily $T(x,1,z)$. It does
not establish that contracting an arbitrary $T(x,y,z)$ with $y>1$
preserves or advances its crossover threshold. Within the audited
subfamily, the exact landscape sweep checked all $134{,}391{,}973$
graphs with $|V| \le 10{,}000{,}000$ against $E_n^{(2)}$ for odd
$n \le 51$ using exact integer arithmetic.

\begin{compresult}[Frontier certification in the tested range]
\label{thm:frontier}
For $T(x,1,z)$, the difference $\Delta_n(x,z) = \Hom(P_n, T) -
\Hom(E_n^{(2)}, T)$ is an exact polynomial in~$x$ and~$z$.
Within the tested range of Table~\ref{tab:landscape_full}, the
minimum-vertex graph achieving first crossover at odd threshold~$n$
is certified by exhaustive exact search over the integer pairs
$(x,z)$ satisfying $\Delta_n(x,z) > 0$ and $\Delta_{n-2}(x,z) \le 0$.
The dominant cubic of $\Delta_n$ explains the geometric shape of the
feasible region but is not, by itself, a complete proof of minimality
at the small thresholds realised here.
\end{compresult}

\begin{proof}[Derivation of the feasible wedge]
Both $\Hom(P_n, T)$ and $\Hom(E_n^{(2)}, T)$ are computed via
finite matrix multiplications over the $4 \times 4$ quotient matrix.
Their difference is therefore an exact polynomial.  For example:
\begin{align*}
  \Delta_7(x,z) &= -x\bigl(x^3 - x^2 - 2xz^2 - xz - x + (z+1)^3\bigr).
\end{align*}
Setting $\Delta_7 > 0$ and extracting the dominant cubic,
$-x^3 + 2xz^2 - z^3 > 0$, the substitution $w = x/z$ yields the
characteristic equation $-w^3 + 2w - 1 = 0$, which factors as
$-(w-1)(w^2 + w - 1) = 0$.  The positive roots give
\[
  \frac{\sqrt{5} - 1}{2} < \frac{x}{z} < 1,
\]
i.e., the dominant terms place the feasible integer pairs inside a
geometric wedge bounded by the inverse golden ratio
$\varphi^{-1} \approx 0.618$. This explains the observed frontier
shape and sharply narrows the exact search region, but the final
minimality claims below are certified by exhaustive integer search in
the tested range rather than by the wedge alone. For example, at
$n = 7$ the
unique minimum is $(x,z) = (16,19)$ with $\Delta_7(16,19) = 512 > 0$
and $\Delta_5(16,19) = -560 \le 0$, yielding $|V| = 337$.  This
improves the $1{,}801$-vertex bound of
Galvin et al.~\cite{galvin2025} to 337 vertices. Exact exhaustive search then certifies the remaining entries in
Table~\ref{tab:landscape_full}: $(9,11)$ at $n = 9$; $(8,10)$ at
$n = 11$; $(7,9)$ at $n = 13$; and $(8,12)$ at $n = 15$, matching
the computational landscape in Table~\ref{tab:landscape_full}.
\end{proof}

\begin{remark}[Algebraic ballast in the strongly Leontovich case]
\label{rem:ballast}
A graph is \emph{strongly Leontovich}~\cite{galvin2026} if
$E_n$ beats the path asymptotically, i.e.\ the leading coefficient ratio
satisfies $\rho_d(H) < 1$ in Eq.~\eqref{eq:rho}.  This is harder than
the standard Leontovich condition, which only requires destructive
interference at \emph{some} odd~$n$.

Under 4-orbit symmetry, the quotient dynamics are governed by a
$4 \times 4$ quotient matrix with only 3 integer degrees of freedom
$(x, y, z)$.  To force the strict ratio inequality $\rho_d(H) < 1$ in
Eq.~\eqref{eq:rho} within this tiny parameter space, the branching degrees
must be scaled to extreme values (e.g., $x = 400$, $z = 800$ for
$\hat{T}(400, 3, 800)$).  The resulting $\sim\!960{,}000$ vertices are
\emph{algebraic ballast}: redundant, identical branches imposed by the
symmetry constraint, carrying eigenvectors orthogonal to the initial state
that project to zero in the homomorphism counts and contribute nothing to
the crossover.

Adding a 5th orbit expands the quotient matrix from $4 \times 4$ to
$5 \times 5$, unlocking a new independent variable.  This produces many
small ordinary Leontovich examples, but the strong condition must be checked
asymptotically rather than by observing a finite crossover window.  A corrected
exact audit of 5-orbit looped symmetric trees
$\hat{T}(d_1, d_2, d_3, d_4)$ finds that the smallest verified
\textit{strongly} Leontovich example in this sweep is
$\hat{T}(1, 35, 1, 50)$ on $1{,}822$ vertices.  This still improves the
$961{,}601$-vertex 4-orbit upper bound by roughly a factor of $528$, while the
smaller examples in Table~\ref{tab:strong_frontier} are ordinary finite-window
Leontovich graphs rather than strongly Leontovich graphs.
For every row in the table, the sign of $\rho_2-1$ is additionally certified
over the exact algebraic Perron root: the verifier isolates that root in a
rational interval and proves that the corresponding integer sign polynomial
has no zero in the interval. Thus these asymptotic classifications do not
depend on a floating-point tolerance.

\begin{table}[H]
\centering
\caption{Audit of 5-orbit looped symmetric trees: ordinary versus strongly Leontovich behavior.}
\label{tab:strong_frontier}
\begin{tabular}{@{}ccll@{}}
\toprule
Opening $n$ & $|V|$ & Parameters $(d_1, d_2, d_3, d_4)$ & Verdict \\
\midrule
7 & 1,066 & $\hat{T}(1, 28, 1, 36)$ & ordinary only; closes at $n=1827$ \\
9 & 370 & $\hat{T}(1, 16, 1, 21)$ & ordinary only; closes at $n=251$ \\
11 & 262 & $\hat{T}(1, 13, 1, 18)$ & ordinary only; closes at $n=141$ \\
13 & 242 & $\hat{T}(1, 12, 1, 18)$ & ordinary only; closes at $n=127$ \\
15 & 546 & $\hat{T}(1, 17, 1, 30)$ & ordinary only; closes at $n=261$ \\
17 & 422 & $\hat{T}(1, 15, 1, 26)$ & ordinary only; closes at $n=133$ \\
\midrule
7 & 1,822 & $\hat{T}(1, 35, 1, 50)$ & strongly Leontovich (smallest known record) \\
\bottomrule
\end{tabular}
\end{table}
\end{remark}

\subsection*{Non-oscillatory threshold}

\begin{proposition}
Suppose the odd-$n$ difference has the exact two-mode form
\[
  \Delta_n = \Hom(P_n, H) - \Hom(E_n^{(d)}, H)
  = A\,\lambda_1^n + B\,\lambda_2^n
\]
for some $A>0$, $B<0$, and $0 < \lambda_2 < \lambda_1$. Then
$\Delta_n$ restricted to odd~$n$ can change sign at most once. This
statement does not apply to the 5-orbit looped symmetric targets in
Table~\ref{tab:strong_frontier}.
\end{proposition}

\begin{proof}[Proof sketch]
Write $r = \lambda_2/\lambda_1 \in (0,1)$. Then
\[
  \Delta_n = \lambda_1^n\bigl(A + B r^n\bigr)
\]
and $B r^n$ is strictly increasing because $B<0$ and $r^n$ decreases
with $n$. Hence $A + B r^n$ crosses zero at most once on the odd
subsequence. Any finite base cases still need to be checked directly.
\end{proof}

\noindent
We verified this computationally for the 76-vertex graph: $\Delta$
is negative at odd $n \le 15$ and strictly positive at odd
$17\le n\le43$, with the gap growing
exponentially from $\Delta_{17} = 1{,}993{,}646$ to
$\Delta_{43} > 10^{17}$.

\subsection*{A Bipartite Parity Hypothesis}

\begin{remark}[Bipartite parity hypothesis (refuted in Section~\ref{sec:doublecover})]
For any bipartite target graph $H$, $\Hom(E_n^{(d)}, H) \ge
\Hom(P_n, H)$ for all even~$n$.  That is, the path $P_n$ is
always the minimizer among trees on $n$ vertices when $n$ is even
and $H$ is bipartite.
\end{remark}

\noindent
Extensive finite searches suggested this hypothesis: across all tested
target graphs (the 608,930,559 connected bipartite graphs on $m \le 15$
vertices, together with all source trees through $n=22$, including
$5{,}623{,}756$ at $n=22$), no even-$n$ crossover was observed.
Section~\ref{sec:doublecover}
disproves it: the bipartite double cover $B$ of $\hat T(1,35,1,50)$
overtakes the path at the even threshold $n = 17{,}340$.

\begin{remark}
The same parity effect appears for independent sets: the consecutive
ratios $a(N)/a(N-1)$ oscillate with period~2; the looped target
$H_{\mathrm{ind}}$ breaks bipartiteness, so the parity protection is
gone.
\end{remark}
